\section{AIME I 1999}
        \begin{enumerate}
        	\item (\textit{AIME I 1999}) Find the smallest prime that is the fifth term of an increasing arithmetic sequence, all four preceding terms also being prime.


 
 \textbf{Answer: }029



\textbf{Solution}

Obviously, all of the terms must be \href{https://artofproblemsolving.com/wiki/index.php?title=Odd}{odd}. The common difference between the terms cannot be $2$ or $4$, since otherwise there would be a number in the sequence that is divisible by $3$. However, if the common difference is $6$, we find that $5,11,17,23$, and $29$ form an \href{https://artofproblemsolving.com/wiki/index.php?title=Arithmetic\_sequence}{arithmetic sequence}. Thus, the answer is $029$.

 \textbf{Alternate Solution} If we let the arithmetic sequence to be $p, p+a, p+2a, p+3a$, and $p+4a$, where $p$ is a prime number and $a$ is a positive integer, we can see that $p$ cannot be multiple of $2$ or $3$ or $4$. Smallest such prime number is $5$, and from a quick observation we can see that when $a$ is $6$, the terms of the sequence are all prime numbers. The sequence becomes $5, 11, 17, 23, 29$, so the answer is $029$.

 \textbf{Additionally} Even without observing these patterns, while rapid guessing and checking this information also appears. We can treat every ones digit of a prime number (which are all odd except for 5) as a value mod 5. The common difference must be even, so visualizing every even number as a value mod 5 can rule out every answer that doesn't include 5 in the sequence. Namely, normally if the ones digit is 5 then it would not be prime.

 
1 mod 5: since there are 5 numbers one will inevitably become divisible by 5.

 2 mod 5: guess and check or noticing 1/3 + first 4 even numbers will give 0 mod 5.

 3 mod 5: same as -2 mod 5.

 4 mod 5: again, same as -1 mod 5.

 
Then we can conclude every sequence will contain a 5 in the ones digit, so we are directed towards 5 and given our answer.

 
-jackshi2006

 




	\item (\textit{AIME I 1999}) Consider the parallelogram with vertices $(10,45),$ $(10,114),$ $(28,153),$ and $(28,84).$  A line through the origin cuts this figure into two congruent polygons.  The slope of the line is $m/n,$ where $m_{}$ and $n_{}$ are relatively prime positive integers.  Find $m+n.$


 
 \textbf{Answer: }118



\textbf{Solution 1}

Let the first point on the line $x=10$ be $(10,45+a)$ where a is the height above $(10,45)$.  Let the second point on the line $x=28$ be $(28, 153-a)$. For two given points, the line will pass the origin if the coordinates are \href{https://artofproblemsolving.com/wiki/index.php?title=Proportion}{proportional} (such that $\frac{y_1}{x_1} = \frac{y_2}{x_2}$). Then, we can write that $\frac{45 + a}{10} = \frac{153 - a}{28}$. Solving for $a$ yields that $1530 - 10a = 1260 + 28a$, so $a=\frac{270}{38}=\frac{135}{19}$. The slope of the line (since it passes through the origin) is $\frac{45 + \frac{135}{19}}{10} = \frac{99}{19}$, and the solution is $m + n = \boxed{118}$.

 





\textbf{Solution 2 (the best solution)}

You can clearly see that a line that cuts a parallelogram into two congruent pieces must go through the center of the parallelogram. Taking the midpoint of $(10,45)$, and $(28,153)$ gives $(19,99)$, which is the center of the parallelogram. Thus the slope of the line must be $\frac{99}{19}$, and the solution is $\boxed{118}$.

 Edit by ngourise: if u dont think this is the best solution then smh

 





\textbf{Solution 3 (Area)}

Note that the area of the parallelogram is equivalent to $69 \cdot 18 = 1242,$ so the area of each of the two trapezoids with congruent area is $621.$ Therefore, since the height is $18,$ the sum of the bases of each trapezoid must be $69.$ 

 The points where the line in question intersects the long side of the parallelogram can be denoted as $(10, \frac{10m}{n})$ and $(28, \frac{28m}{n}),$ respectively. We see that $\frac{10m}{n} - 45 + \frac{28m}{n} - 84 = 69,$ so $\frac{38m}{n} = 198 \implies \frac{m}{n} = \frac{99}{19} \implies \boxed{118}.$

 Solution by Ilikeapos

 





\textbf{Solution 4 (centroid) (fastest way)}

$(\Sigma x_i /4, \Sigma y_i /4)$ is the centroid, which generates $(19,99)$, so the answer is $\boxed{118}$. This is the fastest way because you do not need to find the opposite vertices by drawing.

 Solution by maxamc

 




	\item (\textit{AIME I 1999}) Find the sum of all positive integers $n$ for which $n^2-19n+99$ is a perfect square.


 
 \textbf{Answer: }038



\textbf{Solution 1}

If $n^2-19n+99=x^2$ for some positive integer $x$, then rearranging we get $n^2-19n+99-x^2=0$. Now from the quadratic formula, 

 
            \begin{enumerate}[label=(\roman*)]
                \item $n=\frac{19\pm \sqrt{4x^2-35}}{2}$
            \end{enumerate}
            Because $n$ is an integer, this means $4x^2-35=q^2$ for some nonnegative integer $q$. Rearranging gives $(2x+q)(2x-q)=35$. Thus $(2x+q, 2x-q)=(35, 1)$ or $(7,5)$, giving $x=3$ or $9$. This gives $n=1, 9, 10,$ or $18$, and the sum is $1+9+10+18=\boxed{38}$.

 





\textbf{Solution 2}

Suppose there is some $k$ such that $x^2 - 19x + 99 = k^2$. Completing the square, we have that $(x - 19/2)^2 + 99 - (19/2)^2 = k^2$, that is, $(x - 19/2)^2 + 35/4 = k^2$. Multiplying both sides by 4 and rearranging, we see that $(2k)^2 - (2x - 19)^2 = 35$. Thus, $(2k - 2x + 19)(2k + 2x - 19) = 35$. We then proceed as we did in the previous solution.

 





\textbf{Solution 3}

When $n \geq 12$, we have 
 \[(n-10)^2 < n^2 -19n + 99 < (n-8)^2.\]

 So if $n \geq 12$ and $n^2 -19n + 99$ is a perfect square, then 
 \[n^2 -19n + 99 = (n-9)^2\]

 or $n = 18$.

 For $1 \leq n < 12$, it is easy to check that $n^2 -19n + 99$ is a perfect square when $n = 1, 9$ and $10$ ( using the identity $n^2 -19n + 99 = (n-10)^2 + n - 1.)$

 We conclude that the answer is $1 + 9 + 10 + 18 = \boxed{38}.$

 





\textbf{Solution 4 (Factoring and Midpoint Substitution)}

The problem states that we want to find integers $n$ such that, 
 \[n^2 - 19n + 99 = k^2 \text{, for some integer } k \in \mathbb{Z}\]Notice that the expression can be rewritten to,
 \[n^2 - 19n + 18 = k^2 - 81\]
 \[(n-18)(n-1) = k^2 - 81\]

 Using midpoint substitution, let $u = n - \frac{18+1}{2} = n - \frac{19}{2}$

 Substituting to the factored expression yields,
 \[\left(u-\frac{17}{2}\right)\left(u+\frac{17}{2}\right) = k^2 - 81\]
 \[u^2 - \left(\frac{17}{2}\right)^2 = k^2 - 81\]
 \[u^2 - k^2 = -\frac{35}{4}\]
 \[4(u+k)(u-k) = -35\]
 \[(2u + 2k)(2u - 2k) = -35\]Notice the factors of $-$$35$ are $(-5, 7), (5, -7), (-1, 35)$, and $(1, -35)$

 Solving the systems of equations for each factored pair gives $n = 1, 9, 10, 18$ and $1 + 9 + 10 + 18 = \boxed{38}$

 




	\item (\textit{AIME I 1999}) The two squares shown share the same center $O_{}$ and have sides of length 1. The length of $\overline{AB}$ is $\frac{43}{99}$ and the area of octagon $ABCDEFGH$ is $\frac{m}{n}$ where $m_{}$ and $n_{}$ are relatively prime positive integers.  Find $m+n.$


 
                    \begin{figure}[H]
                        \centering
                        \includegraphics[width=0.5\textwidth]{Problem/AIME_1999_I_4_Problem_Figure_1.png}
                    \end{figure}
                    


 
 \textbf{Answer: }185



\textbf{Solution 1}

Triangles $AOB$, $BOC$, $COD$, etc. are congruent by symmetry (you can prove it rigorously by using the power of a point to argue that exactly two chords of length $1$ in the circumcircle of the squares pass through $B$, etc.), and each area is $\frac{\frac{43}{99}\cdot\frac{1}{2}}{2}$. Since the area of a triangle is $bh/2$, the area of all $8$ of them is $\frac{86}{99}$ and the answer is $\boxed{185}$.

 





\textbf{Solution 2}

Define the two possible \href{https://artofproblemsolving.com/wiki/index.php?title=Distance}{distances} from one of the labeled points and the \href{https://artofproblemsolving.com/wiki/index.php?title=Vertex}{corners} of the square upon which the point lies as $x$ and $y$. The area of the \href{https://artofproblemsolving.com/wiki/index.php?title=Octagon}{octagon} (by \href{https://artofproblemsolving.com/wiki/index.php?title=Subtraction}{subtraction} of areas) is $1 - 4\left(\frac{1}{2}xy\right) = 1 - 2xy$. 

 By the \href{https://artofproblemsolving.com/wiki/index.php?title=Pythagorean\_theorem}{Pythagorean theorem},
 \[x^2 + y^2 = \left(\frac{43}{99}\right)^2\]

 Also,
 \begin{align*}x + y + \frac{43}{99} &= 1\\ x^2 + 2xy + y^2 &= \left(\frac{56}{99}\right)^2\end{align*}


 Substituting,
 \begin{align*}\left(\frac{43}{99}\right)^2 + 2xy &= \left(\frac{56}{99}\right)^2 \\ 2xy = \frac{(56 + 43)(56 - 43)}{99^2} &= \frac{13}{99} \end{align*}


 Thus, the area of the octagon is $1 - \frac{13}{99} = \frac{86}{99}$, so $m + n = \boxed{185}$.

 




	\item (\textit{AIME I 1999}) For any positive integer $x_{}$, let $S(x)$ be the sum of the digits of $x_{}$, and let $T(x)$ be $|S(x+2)-S(x)|.$  For example, $T(199)=|S(201)-S(199)|=|3-19|=16.$  How many values of $T(x)$ do not exceed 1999?


 
 \textbf{Answer: }223



\textbf{Solution}

For most values of $x$, $T(x)$ will equal $2$. For those that don't, the difference must be bumping the number up a ten, a hundred, etc. If we take $T(a999)$ as an example,
 \[|(a + 1) + 0 + 0 + 1 - (a + 9 + 9 + 9)| = |2 - 9(3)|\]And in general, the values of $T(x)$ will then be in the form of $|2 - 9n| = 9n - 2$. From $7$ to $1999$, there are $\left\lceil \frac{1999 - 7}{9}\right\rceil = 222$ solutions; including $2$ and there are a total of $\boxed{223}$ solutions.

 




	\item (\textit{AIME I 1999}) A transformation of the first quadrant of the coordinate plane maps each point $(x,y)$ to the point $(\sqrt{x},\sqrt{y}).$  The vertices of quadrilateral $ABCD$ are $A=(900,300), B=(1800,600), C=(600,1800),$ and $D=(300,900).$  Let $k_{}$ be the area of the region enclosed by the image of quadrilateral $ABCD.$  Find the greatest integer that does not exceed $k_{}.$


 
 \textbf{Answer: }314



\textbf{Solution}


 \begin{eqnarray*}A' = & (\sqrt {900}, \sqrt {300})\\ B' = & (\sqrt {1800}, \sqrt {600})\\ C' = & (\sqrt {600}, \sqrt {1800})\\ D' = & (\sqrt {300}, \sqrt {900}) \end{eqnarray*}

 First we see that lines passing through $AB$ and $CD$ have \href{https://artofproblemsolving.com/wiki/index.php?title=Equation}{equations} $y = \frac {1}{3}x$ and $y = 3x$, respectively. Looking at the points above, we see the equations for $A'B'$ and $C'D'$ are $y^2 = \frac {1}{3}x^2$ and $y^2 = 3x^2$, or, after manipulation $y = \frac {x}{\sqrt {3}}$ and $y = \sqrt {3}x$, respectively, which are still linear functions. Basically the square of the image points gives back the original points and we could plug them back into the original equation to get the equation of the image lines. 

 Now take a look at $BC$ and $AD$, which have the equations $y = - x + 2400$ and $y = - x + 1200$. The image equations hence are $x^2 + y^2 = 2400$ and $x^2 + y^2 = 1200$, respectively, which are the equations for \href{https://artofproblemsolving.com/wiki/index.php?title=Circle}{circles}. 

 
                    \begin{figure}[H]
                        \centering
                        \includegraphics[width=0.5\textwidth]{Solution/AIME_1999_I_6_Solution_Figure_1.png}
                    \end{figure}
                    

 To find the area between the circles (actually, parts of the circles), we need to figure out the \href{https://artofproblemsolving.com/wiki/index.php?title=Angle}{angle} of the \href{https://artofproblemsolving.com/wiki/index.php?title=Arc}{arc}. This could be done by $\arctan \sqrt {3} - \arctan \frac {1}{\sqrt {3}} = 60^\circ - 30^\circ = 30^\circ$. So the requested areas are the area of the enclosed part of the smaller circle subtracted from the area enclosed by the part of the larger circle = $\frac {30^\circ}{360^\circ}(R^2\pi - r^2\pi) = \frac {1}{12}(2400\pi - 1200\pi) = 100\pi$. Hence the answer is $\boxed{314}$.

 




	\item (\textit{AIME I 1999}) There is a set of 1000 switches, each of which has four positions, called $A, B, C$, and $D$.  When the position of any switch changes, it is only from $A$ to $B$, from $B$ to $C$, from $C$ to $D$, or from $D$ to $A$.  Initially each switch is in position $A$.  The switches are labeled with the 1000 different integers $(2^{x})(3^{y})(5^{z})$, where $x, y$, and $z$ take on the values $0, 1, \ldots, 9$.  At step $i$ of a 1000-step process, the $i$-th switch is advanced one step, and so are all the other switches whose labels divide the label on the $i$-th switch.  After step 1000 has been completed, how many switches will be in position $A$?


 
 \textbf{Answer: }650



\textbf{Solution}

For each $i$th switch (designated by $x_{i},y_{i},z_{i}$), it advances \textit{itself} only one time at the $i$th step; thereafter, only a switch with larger $x_{j},y_{j},z_{j}$ values will advance the $i$th switch by one step provided $d_{i}= 2^{x_{i}}3^{y_{i}}5^{z_{i}}$ divides $d_{j}= 2^{x_{j}}3^{y_{j}}5^{z_{j}}$. Let $N = 2^{9}3^{9}5^{9}$ be the max switch label. To find the divisor multiples in the range of $d_{i}$ to $N$, we consider the exponents of the number $\frac{N}{d_{i}}= 2^{9-x_{i}}3^{9-y_{i}}5^{9-z_{i}}$. In general, the divisor-count of $\frac{N}{d}$ must be a multiple of 4 to ensure that a switch is in position A:

 $4n = [(9-x)+1] [(9-y)+1] [(9-z)+1] = (10-x)(10-y)(10-z)$, where $0 \le x,y,z \le 9.$We consider the cases where the 3 factors above do not contribute multiples of 4.

 
                \begin{itemize}
                    \item Case of no 2's:
                \end{itemize}


                
            \begin{enumerate}[label=(\roman*)]
                \item The switches must be $(\mathrm{odd})(\mathrm{odd})(\mathrm{odd})$. There are $5$ \href{https://artofproblemsolving.com/wiki/index.php?title=Odd\_integer}{odd integers} in $0$ to $9$, so we have $5 \times 5 \times 5 = 125$ ways.
            \end{enumerate}
            
                \begin{itemize}
                    \item Case of a single 2:
                \end{itemize}


                
            \begin{enumerate}[label=(\roman*)]
                \item The switches must be one of $(2\cdot \mathrm{odd})(\mathrm{odd})(\mathrm{odd})$ or $(\mathrm{odd})(2 \cdot \mathrm{odd})(\mathrm{odd})$ or $(\mathrm{odd})(\mathrm{odd})(2 \cdot \mathrm{odd})$.
            \end{enumerate}
            
            \begin{enumerate}[label=(\roman*)]
                \item Since $0 \le x,y,z \le 9,$ the terms $2\cdot 1, 2 \cdot 3,$ and $2 \cdot 5$ are three valid choices for the $(2 \cdot odd)$ factor above.
            \end{enumerate}
            
            \begin{enumerate}[label=(\roman*)]
                \item We have ${3\choose{1}} \cdot 3 \cdot 5^{2}= 225$ ways.
            \end{enumerate}
            The number of switches in position A is $1000-125-225 = \boxed{650}$.

 




	\item (\textit{AIME I 1999}) Let $\mathcal{T}$ be the set of ordered triples $(x,y,z)$ of nonnegative real numbers that lie in the plane $x+y+z=1.$  Let us say that $(x,y,z)$ supports $(a,b,c)$ when exactly two of the following are true: $x\ge a, y\ge b, z\ge c.$  Let $\mathcal{S}$ consist of those triples in $\mathcal{T}$ that support $\left(\frac 12,\frac 13,\frac 16\right).$  The area of $\mathcal{S}$ divided by the area of $\mathcal{T}$ is $m/n,$  where $m_{}$ and $n_{}$ are relatively prime positive integers, find $m+n.$


 
 \textbf{Answer: }025



\textbf{Solution}

This problem just requires a good diagram and strong 3D visualization.

 
                    \begin{figure}[H]
                        \centering
                        \includegraphics[width=0.5\textwidth]{Solution/AIME_1999_I_8_Solution_Figure_1.png}
                    \end{figure}
                    

 The region in $(x,y,z)$ where $x \ge \frac{1}{2}, y \ge \frac{1}{3}$ is that of a little triangle on the bottom of the above diagram, of $y \ge \frac{1}{3}, z \ge \frac{1}{6}$ is the triangle at the right, and $x \ge \frac 12, z \ge \frac 16$ the triangle on the left, where the triangles are coplanar with the large equilateral triangle formed by $x+y+z=1,\ x,y,z \ge 0$. We can check that each of the three regions mentioned fall under exactly two of the inequalities and not the third. 

 
                    \begin{figure}[H]
                        \centering
                        \includegraphics[width=0.5\textwidth]{Solution/AIME_1999_I_8_Solution_Figure_1.png}
                    \end{figure}
                    

 The side length of the large equilateral triangle is $\sqrt{2}$, which we can find using 45-45-90 $\triangle$ with the axes. Using the formula $A = \frac{s^2\sqrt{3}}{4}$ for \href{https://artofproblemsolving.com/wiki/index.php?title=Equilateral\_triangle}{equilateral triangles}, the area of the large triangle is $\frac{(\sqrt{2})^2\sqrt{3}}{4} = \frac{\sqrt{3}}{2}$. Since the lines of the smaller triangles are \href{https://artofproblemsolving.com/wiki/index.php?title=Parallel}{parallel} to those of the large triangle, by corresponding angles we see that all of the triangles are \href{https://artofproblemsolving.com/wiki/index.php?title=Similar\_triangles}{similar}, so they are all equilateral triangles. We can solve for their side lengths easily by subtraction, and we get $\frac{\sqrt{2}}{6}, \frac{\sqrt{2}}{3}, \frac{\sqrt{2}}{2}$. Calculating their areas, we get $\frac{\sqrt{3}}{8}, \frac{\sqrt{3}}{18}, \frac{\sqrt{3}}{72}$. The \href{https://artofproblemsolving.com/wiki/index.php?title=Ratio}{ratio} $\frac{\mathcal{S}}{\mathcal{T}} = \frac{\frac{9\sqrt{3} + 4\sqrt{3} + \sqrt{3}}{72}}{\frac{\sqrt{3}}{2}} = \frac{14}{36} = \frac{7}{18}$, and the answer is $m + n = \boxed{025}$.

 To simplify the problem, we could used the fact that the area ratios are equal to the side ratios squared, and we get $\left(\frac{1}{2}\right)^2 + \left(\frac{1}{3}\right)^2 + \left(\frac{1}{6}\right)^2 = \frac{14}{36} = \frac{7}{18}$.

 




	\item (\textit{AIME I 1999}) A function $f$ is defined on the complex numbers by $f(z)=(a+bi)z,$ where $a_{}$ and $b_{}$ are positive numbers.  This function has the property that the image of each point in the complex plane is equidistant from that point and the origin.  Given that $|a+bi|=8$ and that $b^2=m/n,$ where $m_{}$ and $n_{}$ are relatively prime positive integers.  Find $m+n.$


 
 \textbf{Answer: }259



\textbf{Solution 1}

Suppose we pick an arbitrary point on the \href{https://artofproblemsolving.com/wiki/index.php?title=Complex\_plane}{complex plane}, say $(1,1)$. According to the definition of $f(z)$,  
 \[f(1+i) = (a+bi)(1+i) = (a-b) + (a+b)i,\] this image must be equidistant to $(1,1)$ and $(0,0)$. Thus the image must lie on the line with slope $-1$ and which passes through $\left(\frac 12, \frac12\right)$, so its graph is $x + y = 1$. Substituting $x = (a-b)$ and $y = (a+b)$, we get $2a = 1 \Rightarrow a = \frac 12$. 

 By the \href{https://artofproblemsolving.com/wiki/index.php?title=Pythagorean\_Theorem}{Pythagorean Theorem}, we have $\left(\frac{1}{2}\right)^2 + b^2 = 8^2 \Longrightarrow b^2 = \frac{255}{4}$, and the answer is $\boxed{259}$.

 





\textbf{Solution 2}

Plugging in $z=1$ yields $f(1) = a+bi$. This implies that $a+bi$ must fall on the line $Re(z)=a=\frac{1}{2}$, given the equidistant rule. By $|a+bi|=8$, we get $a^2 + b^2 = 64$, and plugging in $a=\frac{1}{2}$ yields $b^2=\frac{255}{4}$. The answer is thus $\boxed{259}$.

 





\textbf{Solution 3}

We are given that $(a + bi)z$ is equidistant from the origin and $z.$  This translates to
 \begin{eqnarray*} |(a + bi)z - z| & = & |(a + bi)z| \\ |z(a - 1 + bi)| & = & |z(a + bi)| \\ |z|\cdot|(a - 1) + bi| & = & |z|\cdot|a + bi| \\ |(a - 1) + bi| & = & |a + bi| \\ (a - 1)^2 + b^2 & = & a^2 + b^2 \\ & \Rightarrow & a = \frac 12 \end{eqnarray*}Since $|a + bi| = 8,$ $a^2 + b^2 = 64.$ Because $a = \frac 12,$ thus $b^2 = \frac {255}4.$  So the answer is $\boxed{259}$.

 





\textbf{Solution 4}

Let $P$ and $Q$ be the points in the complex plane represented by $z$ and $(a+bi)z$, respectively.  $|a+bi| = 8$ implies $OQ = 8OP$.  Also, we are given $OQ = PQ$, so $OPQ$ is isosceles with base $OP$.  Notice that the base angle of this isosceles triangle is equal to the argument $\theta$ of the complex number $a + bi$, because $(a+bi)z$ forms an angle of $\theta$ with $z$.  Drop the altitude/median from $Q$ to base $OP$, and you end up with a right triangle that shows $\cos \theta = \frac{\frac{1}{2}OP}{8OQ} = \frac{\frac{1}{2}|z|}{8|z|} = \frac{1}{16}$.  Since $a$ and $b$ are positive, $z$ lies in the first quadrant and $\theta < \pi/2$; hence by right triangle trigonometry $\sin \theta = \frac{\sqrt{255}}{16}$.  Finally, $b = |a+bi|\sin\theta = 8\frac{\sqrt{255}}{16} = \frac{\sqrt{255}}{2}$, and $b^2 = \frac{255}{4}$, so the answer is $259$.

 





\textbf{Solution 5}

Similarly to in Solution 3, we see that $|(a + bi)z - z| = |(a + bi)z|$. Letting the point $z = c + di$, we have $\sqrt{(ab+bc-d)^2+(ac-bd-c)^2} = \sqrt{(ac-bd)^2+(ad+bc)^2}$. Expanding both sides of this equation (after squaring, of course) and canceling terms, we get $(d^2+c^2)(-2a+1) = 0$. Of course, $(d^2+c^2)$ can't be zero because this property of the function holds for all complex $z$. Therefore, $a = \frac{1}{2}$ and we proceed as above to get $\boxed{259}$.

 





\textbf{Solution 6}

This is a solution that minimizes the use of complex numbers, turning this into an introductory algebra analytic geometry problem.

 Consider any complex number $z=c+di$. Let $z$ denote point $P$ on the complex plane. Then $P=(c,d)$ on the complex plane. The equation for the line $OP$ is $y=\frac{d}{c}x$.

 Let the image of point $P$ be $Q$, after the point undergoes the function. Since each image is equidistant from the preimage and the origin, $Q$ must be on the perpendicular bisector of $OP$.Given $z=c+di$, $f(z)=(a+bi)(c+di)=(ac-bd)+(ad+bc)i$. Then $Q=(ac-bd,ad+bc)$. The midpoint of $OP$ is $(0.5c, 0.5d)$. Since the slopes of two respectively nonvertical and nonhorizontal lines have a product of $-1$, using the point-slope form, the equation of the perpendicular line to $OP$ is $y-0.5d=-\frac{c}{d}(x-0.5c)$. Rearranging, we have $y=-\frac{cx}{d}+\frac{c^2}{2d}+\frac{d}{2}$.

 Since we know that $Q=(ac-bd,ad+bc)$, thus we plug in $Q$ into the line: $ad+bc=-\frac{ac^2-bcd}{d}+\frac{c^2}{2d}+\frac{d}{2}$.

 Let's start canceling. $2ad^2+2bcd=-2ac^2+2bcd+c^2+d^2$. Subtracting, $c^2+d^2-2ac^2=2ad^2$. Thus $c^2+d^2=2ac^2+2ad^2$. Since this is an identity for any $(c,d)$, thus $2a=1$. $a=\frac{1}{2}$. Since $|a+bi|=8$, thus $a^2+b^2=64$ (or simply think of $a+bi$ as the point $(a,b)$, and $|a+bi|$ being the distance of $(a,b)$ to the origin). Thus plug in $a=\frac{1}{2}, b^2=\frac{255}{4}$. Since $255$ and $4$ are relatively prime, the final result is $255+4=\boxed{259}$.

 




	\item (\textit{AIME I 1999}) Ten points in the plane are given, with no three collinear.  Four distinct segments joining pairs of these points are chosen at random, all such segments being equally likely.  The probability that some three of the segments form a triangle whose vertices are among the ten given points is $m/n,$ where $m_{}$ and $n_{}$ are relatively prime positive integers.  Find $m+n.$


 
 \textbf{Answer: }489



\textbf{Solution 1}

First, let us find the number of triangles that can be formed from the 10 points. Since none of the points are collinear, it is possible to pick ${10\choose3}$ sets of 3 points which form triangles. However, a fourth distinct segment must also be picked. Since the triangle accounts for 3 segments, there are $45 - 3 = 42$ segments remaining.

 The total number of ways of picking four distinct segments is ${45\choose4}$. Thus, the requested probability is $\frac{{10\choose3} \cdot 42}{{45\choose4}} = \frac{10 \cdot 9 \cdot 8 \cdot 42 \cdot 4!}{45 \cdot 44 \cdot 43 \cdot 42 \cdot 3!} = \frac{16}{473}$. The solution is $m + n = 489$.

 





\textbf{Solution 2}

Note that 4 segments can NEVER form 2 triangles. Therefore, we just need to multiply the probability that the first three segments picked form a triangle by 4. We can pick any segment for the first choice, then only segments that share an endpoint with the first one, then the one segment that completes the triangle. Note that the fourth segment doesn't matter in this case.Note that there are $(9 - 1) \times 2 = 16$ segments that share an endpoint with the first segment. The answer is then $4 \times \frac{16}{44} \times \frac{1}{43} = \frac{16}{11} \times \frac{1}{43} = \frac{16}{473} \implies m + n = \boxed{489}$-whatRthose

 





\textbf{Solution 3}

Instead of working with the four segments, let's focus on their endpoints. When we select these segments, we are working with $4, 5, 6, 7,$ or $8$ endpoints in total. 

 If we have $6, 7,$ or $8$ endpoints, it is easy to see that we cannot form a triangle by drawing four segments between them, because at least one point will be "left out".

 However, if we have $5$ endpoints, we can use three segments form a triangle with three of the points and connect the remaining two points with the last segment. There are ${10\choose5}$ ways to select these $5$ points from the original $10,$ and ${5\choose3}$ ways to decide which three points are in the triangle.

 Finally, if we have $4$ endpoints, we can also form a triangle with three of the points, then use the remaining segment to connect the last point to either of the previous three. We have ${10\choose4}$ ways to select the $4$ points and ${4\choose3}$ ways to choose three points for the triangle. Finally, we must connect the last point to one vertex of the triangle; we can do this in $3$ ways.

 As in Solution 1, there are ${45\choose4}$ total ways to select four segments. So, our desired probability is

 
 \[\dfrac{{10\choose5}{5\choose3}+{10\choose4}{4\choose3}\cdot 3}{{45\choose4}}=\dfrac{16}{473} \implies m + n = \boxed{489}.\]

 --vaporwave

 




	\item (\textit{AIME I 1999}) Given that $\sum_{k=1}^{35}\sin 5k=\tan \frac mn,$ where angles are measured in degrees, and $m_{}$ and $n_{}$ are relatively prime positive integers that satisfy $\frac mn<90,$ find $m+n.$


 
 \textbf{Answer: }177



\textbf{Solution 1}

Let $s = \sum_{k=1}^{35}\sin 5k  = \sin 5 + \sin 10 + \ldots + \sin 175$. We could try to manipulate this sum by wrapping the terms around (since the first half is equal to the second half), but it quickly becomes apparent that this way is difficult to pull off. Instead, we look to \href{https://artofproblemsolving.com/wiki/index.php?title=Telescope}{telescope} the sum. Using the \href{https://artofproblemsolving.com/wiki/index.php?title=Trigonometric\_identity}{identity} $\sin a \sin b = \frac 12(\cos (a-b) - \cos (a+b))$, we can rewrite $s$ as 

 
 \begin{align*} s \cdot \sin 5 = \sum_{k=1}^{35} \sin 5k \sin 5 &= \sum_{k=1}^{35} \frac{1}{2}(\cos (5k - 5)- \cos (5k + 5))\\ &= 0.5(\cos 0 - \cos 10 + \cos 5 - \cos 15 + \cos 10 \ldots + \cos 165 - \cos 175+ \cos 170 - \cos 180) \end{align*}


 This telescopes to 
 \[s = \frac{\cos 0 + \cos 5 - \cos 175 - \cos 180}{2 \sin 5} = \frac{1 + \cos 5}{\sin 5}.\] Manipulating this to use the identity $\tan x = \frac{1 - \cos 2x}{\sin 2x}$, we get 
 \[s = \frac{1 - \cos 175}{\sin 175} \Longrightarrow s = \tan \frac{175}{2},\] and our answer is $\boxed{177}$.

 





\textbf{Solution 2}

This solution will work entirely in degrees. Let

 
 \[z = e^{5^\circ i}.\]

 For each integer $k$, consider the unit vector in the complex plane
 \[v_k = z^k = e^{5k^\circ i}.\]

 Geometrically, $v_k$ is a unit vector making an angle of $5k ^\circ$ with the positive real axis. Now draw each vector $v_k$ for $1 \leq k \leq 35$ tip-to-tail on the complex plane. Begin at $O$, the origin, and draw the vector $v_1$. Then starting at where $v_1$ ended, draw $v_2$. Continue until you draw $v_{35}$, and the sum of these $35$ vectors will end at $B$, as seen in the diagram below. Because each successive vector is obtained by rotating the previous vector by $5^\circ$, we have a portion of a regular polygon with an exterior angle of $5$ degrees, so it has $\frac{360^\circ}{5^\circ} = 72$ sides. 

 Additionally, notice that $B$ is directly above $O$. This is because the vector sum will have its horizontal components completely cancel out. We realize that $\cos{5} + \cos{175} = 0$, and we can form a similar pair for each of the other complex numbers, except for $v_18 = i$. 

 
                    \begin{figure}[H]
                        \centering
                        \includegraphics[width=0.5\textwidth]{Solution/AIME_1999_I_11_Solution_Figure_1.png}
                    \end{figure}
                    The above diagram starts at the origin $O$ and draws the $35$ vectors. 

 Note that
 \[\text{Im}(v_k) = \sin{5k}.\] Therefore, $\sum_{k=1}^{35}\sin 5k$ is simply the y-coordinate of the vector sum $v_1 + v_2 + \dots v_{35}$. In the diagram, this is the distance of the vertical blue line connecting $O$ to the topmost point.

 Let $A$ be the vertex of the $72$-gon that is directly to the left of $O$ and $B$ be the resulting vertex of summing the $35$ vectors. Looking at right triangle $AOB$, we notice that $AO=1$. Remember that our desired sum equals $BO = \tan{\angle BAO}.$ We now just need to find $\angle BAO$.

 Notice that $BA$ bisects $A$'s interior angle, so $\angle BAO$ is half of the interior angle of the $72$-gon. So $\angle BAO = \frac{175}{2}$.

 Recall that our sum equals $\tan{\frac{m}{n}} = \tan{\frac{175}{2}}$, so our answer is $175 + 2 = \boxed{177}$.

 





\textbf{Solution 3}

We note that $\sin x = \mbox{Im } e^{ix}\text{*}$. We thus have that
 \begin{align*} \sum_{k = 1}^{35} \sin 5k &= \sum_{k = 1}^{35} \mbox{Im } e^{5ki}\\ &= \mbox{Im } \sum_{k = 1}^{35} e^{5ki}\\ &= \mbox{Im } \frac{e^{5i}(1 - e^{180i})}{1 - e^{5i}}\\ &= \mbox{Im } \frac{2\cos5 + 2i \sin 5}{(1 - \cos 5) - i \sin 5}\\ &= \mbox{Im } \frac{(2 \cos 5 + 2i \sin 5)[(1 - \cos 5) + i \sin 5]}{(1 - \cos 5)^2 + \sin^2 5}\\ &= \frac{2 \sin 5}{2 - 2 \cos 5}\\ &= \frac{\sin 5}{1 - \cos 5}\\ &= \frac{\sin 175}{1 + \cos 175} \\ &= \tan \frac{175}{2}.\\ \end{align*}
The desired answer is thus $175 + 2 = \boxed{177}$.

 
                \begin{itemize}
                    \item Only if $x$ is in radians, which it is not. However, the solution is still viable, so keep reading.
                \end{itemize}


                





\textbf{Solution 4}

Let $x=e^{\frac{i\pi}{36}}$. By Euler's Formula, $\sin{5k^\circ}=\frac{x^k-\frac{1}{x^{k}}}{2i}$.

 The sum we want is thus $\frac{x-\frac{1}{x}}{2i}+\frac{x^2-\frac{1}{x^{2}}}{2i}+\cdots+\frac{x^{35}-\frac{1}{x^{35}}}{2i}$

 We factor the $\frac{1}{2i}$ and split into two geometric series to get $\frac{1}{2i}\left(\frac{-\frac{1}{x^{35}}(x^{35}-1)}{x-1}+\frac{x(x^{35}-1)}{x-1}\right)$

 However, we note that $x^{36}=-1$, so $-\frac{1}{x^{35}}=x$, so our two geometric series are actually the same. We combine the terms and simplify to get $\frac{1}{i}\left(\frac{x^{36}-x}{x-1}\right)$

 Apply Euler's identity and simplify again to get $\frac{1}{i}\left(\frac{-x-1}{x-1}\right)$

 Now, we need to figure out how to express this as the tangent of something. We note that $\tan(5k^\circ)=\frac{\sin(5k^\circ)}{\cos(5k^\circ)}=\frac{\frac{x^k-\frac{1}{x^k}}{2i}}{\frac{x^k+\frac{1}{x^k}}{2}}=\frac{1}{i}\frac{x^{2k}-1}{x^{2k}+1}$.

 So, we set the two equal to each other to solve for $k$. Cross multiplying gets $(-x-1)(x^{2k}+1)=(x-1)(x^{2k}-1)$. Expanding yields $-x^{2k+1}-x-x^{2k}-1=x^{2k+1}-x-x^{2k}+1$. Simplifying yields $x^{2k+1}=-1$. Since $2k+1=36$ is the smallest solution, we have $k=\frac{35}{2}$, and the argument of tangent is $5k=\frac{175}{2}$. The requested sum is $175+2=\boxed{177}$.

 





\textbf{Solution 5}

This solution is inspired by \href{https://artofproblemsolving.com/wiki/index.php?title=1997\_AIME\_Problems/Problem\_11}{1997 AIME Problems/Problem 11}.First of all, $\sum_{k=1}^{35}\sin 5k=\sum_{k=0}^{35}\sin 5k$ because the difference is $\sin 0 = 0$. Now consider the sum $\sum_{k=0}^{35}\cos 5k$. Since $\cos 5=-\cos 175$, $\cos 10=-\cos 170$ and so on, this simplifies to $\cos 0=1$.Now, consider the ratio
 \[\dfrac{\sum\limits_{n=0}^{35}\sin 5k}{\sum\limits_{n=0}^{35}\cos 5k.}\]Using the sum-to-product identities, the numerator can be written as 
 \[\sin \frac{175}{2} (\cos \frac{175}{2} + \cos \frac {165}{2}+\cdots + \cos \frac{5}{2}.)\]Similarly (via sum-to-product identities), the denominator is equivalent to 
 \[\cos \frac{175}{2} (\cos \frac{175}{2} + \cos \frac {165}{2}+\cdots + \cos \frac{5}{2}.)\]If we substitute these equivalent expressions into the original ratio, many parts cancel. We are left with
 \[\frac{\sin \frac{175}{2}}{\cos \frac{175}{2}} = \tan{\frac{175}{2}}.\]Aha! Remembering that the original denominator $\sum_{k=0}^{35}\cos 5k$ was equal to 1, we realize that the numerator, $\sum_{k=0}^{35}\sin 5k=\tan{\frac{175}{2}}$. Since $\sum_{k=0}^{35}\sin 5k=\sum_{k=1}^{35}\sin 5k$, we get $\tan \frac{175}{2}$, or $m+n=\boxed{177}$, as our final answer. $\blacksquare$~ewei12

 




	\item (\textit{AIME I 1999}) The inscribed circle of triangle $ABC$ is tangent to $\overline{AB}$ at $P_{},$  and its radius is 21.  Given that $AP=23$ and $PB=27,$ find the perimeter of the triangle.


 
 \textbf{Answer: }345



\textbf{Solution}


                    \begin{figure}[H]
                        \centering
                        \includegraphics[width=0.5\textwidth]{Solution/AIME_1999_I_12_Solution_Figure_1.png}
                    \end{figure}
                    





\textbf{Solution 1}

Let $Q$ be the tangency point on $\overline{AC}$, and $R$ on $\overline{BC}$. By the \href{https://artofproblemsolving.com/wiki/index.php?title=Two\_Tangent\_Theorem}{Two Tangent Theorem}, $AP = AQ = 23$, $BP = BR = 27$, and $CQ = CR = x$. Using $rs = A$, where $s = \frac{27 \cdot 2 + 23 \cdot 2 + x \cdot 2}{2} = 50 + x$, we get $(21)(50 + x) = A$. By \href{https://artofproblemsolving.com/wiki/index.php?title=Heron\%27s\_formula}{Heron's formula}, $A = \sqrt{s(s-a)(s-b)(s-c)} = \sqrt{(50+x)(x)(23)(27)}$. Equating and squaring both sides, 

 
 \begin{eqnarray*}  [21(50+x)]^2 &=& (50+x)(x)(621)\\ 441(50+x) &=& 621x\\ 180x = 441 \cdot 50 &\Longrightarrow & x = \frac{245}{2}   \end{eqnarray*}

 We want the perimeter, which is $2s = 2\left(50 + \frac{245}{2}\right) = \boxed{345}$.

 





\textbf{Solution 2}

Let the incenter be denoted $I$.  It is commonly known that the incenter is the intersection of the angle bisectors of a triangle.  So let $\angle ABI = \angle CBI = \alpha, \angle BAI = \angle CAI = \beta,$ and $\angle BCI = \angle ACI = \gamma.$

 We have that
 \begin{eqnarray*} \tan \alpha & = & \frac {21}{27} \\ \tan \beta & = & \frac {21}{23} \\ \tan \gamma & = & \frac {21}x. \end{eqnarray*}So naturally we look at $\tan \gamma.$  But since $\gamma = \frac \pi2 - (\beta + \alpha)$ we have
 \begin{eqnarray*} \tan \gamma & = & \tan\left(\frac \pi2 - (\beta + \alpha)\right) \\ & = & \frac 1{\tan(\alpha + \beta)} \\ \Rightarrow \frac {21}x & = & \frac {1 - \frac {21\cdot 21}{23\cdot 27}}{\frac {21}{27} + \frac {21}{23}} \end{eqnarray*}Doing the algebra, we get $x = \frac {245}2.$

 The perimeter is therefore $2\cdot\frac {245}2 + 2\cdot 23 + 2\cdot 27 = \boxed{345}.$

 


 





\textbf{Solution 3}

Let unknown side has length as $x$, Assume three sides of triangles are $a,b,c$, the area of the triangle is $S$.

 Note that $r=\frac{2S}{a+b+c}=21,S=1050+21x$

 $\tan\angle{\frac{B}{2}}=\frac{7}{9}, \tan\angle{B}=\frac{63}{16}$. Use trig identity, knowing that $1+\cot^2\angle{B}=\csc^2\angle{B}$, getting that $\sin\angle{B}=\frac{63}{65}$

 Now equation $(x+27)*50*\frac{63}{65}*\frac{1}{2}=1050+21x; x=\frac{245}{2}$, the final answer is $245+100=345$~bluesoul

 





\textbf{Solution 4 (Trig Bash)}


                    \begin{figure}[H]
                        \centering
                        \includegraphics[width=0.5\textwidth]{Solution/AIME_1999_I_12_Solution_Figure_1.png}
                    \end{figure}
                    Let \(CR = CQ = x\). Using the inradius formula for the area of a triangle, we have
 \[[ABC] = 1050 + 21x,\]where \([ABC]\) denotes the area of \(\triangle ABC\).

 Next, we find \(\sin \angle CBA\). Note that
 \[\angle CBA = \angle RBP = 2\angle OBP.\]

 First, we find the hypotenuse \(OB\) using the Pythagorean Theorem:
 \[OB^2 = 27^2 + 21^2 = 1170 \quad \Rightarrow \quad OB = 3\sqrt{130}.\]

 Let \(\angle OBP = \theta\). Then
 \[\sin \theta = \frac{21}{3\sqrt{130}}, \qquad \cos \theta = \frac{27}{3\sqrt{130}}.\]

 Using the double-angle identity,
 \[\sin(2\theta) = 2\sin \theta \cos \theta= 2\left(\frac{21}{3\sqrt{130}}\right)\left(\frac{27}{3\sqrt{130}}\right)= \frac{1134}{1170}= \frac{63}{65}.\]

 Now we compute the area of \(\triangle ABC\) using the formula \(\frac{1}{2}ab\sin C\):
 \[[ABC] = \frac{1}{2}(50)(27+x)\left(\frac{63}{65}\right)= \frac{315}{13}(27+x).\]

 Equating the two expressions for the area,
 \[1050 + 21x = \frac{315}{13}(27+x).\]

 Solving,
 \[\frac{13}{15}x + \frac{130}{3} = 27 + x,\]
 \[\frac{49}{3} = \frac{2}{15}x,\]so
 \[x = \frac{245}{2}.\]

 Finally, we compute the perimeter of \(\triangle ABC\):
 \[2(23+27) + 2\left(\frac{245}{2}\right) = \boxed{345}.\]

 




	\item (\textit{AIME I 1999}) Forty teams play a tournament in which every team plays every other ($39$ different opponents) team exactly once.  No ties occur, and each team has a $50 \%$ chance of winning any game it plays.  The probability that no two teams win the same number of games is $m/n,$ where $m_{}$ and $n_{}$ are relatively prime positive integers.  Find $\log_2 n.$


 
 \textbf{Answer: }742



\textbf{Solution}

There are ${40 \choose 2} = 780$ total pairings of teams, and thus $2^{780}$ possible outcomes. In order for no two teams to win the same number of games, they must each win a different number of games. Since the minimum and maximum possible number of games won are 0 and 39 respectively, and there are 40 teams in total, each team corresponds uniquely with some $k$, with $0 \leq k \leq 39$, where $k$ represents the number of games the team won. With this in mind, we see that there are a total of $40!$ outcomes in which no two teams win the same number of games. Further, note that these are all the valid combinations, as the team with 1 win must beat the team with 0 wins, the team with 2 wins must beat the teams with 1 and 0 wins, and so on; thus, this uniquely defines a combination.

 
The desired probability is thus $\frac{40!}{2^{780}}$. We wish to simplify this into the form $\frac{m}{n}$, where $m$ and $n$ are relatively prime. The only necessary step is to factor out all the powers of 2 from $40!$; the remaining number is clearly relatively prime to all powers of 2. 

 
The number of powers of 2 in $40!$ is $\left \lfloor \frac{40}{2} \right \rfloor + \left \lfloor \frac{40}{4} \right \rfloor + \left \lfloor \frac{40}{8} \right \rfloor + \left \lfloor \frac{40}{16} \right \rfloor + \left \lfloor \frac{40}{32} \right \rfloor = 20 + 10 + 5 + 2 + 1 = 38.$ 

 
$780-38 = \boxed{742}$.

 




	\item (\textit{AIME I 1999}) Point $P_{}$ is located inside triangle $ABC$ so that angles $PAB, PBC,$ and $PCA$ are all congruent.  The sides of the triangle have lengths $AB=13, BC=14,$ and $CA=15,$ and the tangent of angle $PAB$ is $m/n,$ where $m_{}$ and $n_{}$ are relatively prime positive integers.  Find $m+n.$


 
 \textbf{Answer: }463



\textbf{Solution}


                    \begin{figure}[H]
                        \centering
                        \includegraphics[width=0.5\textwidth]{Solution/AIME_1999_I_14_Solution_Figure_1.png}
                    \end{figure}
                    





\textbf{Solution 1}

Drop \href{https://artofproblemsolving.com/wiki/index.php?title=Perpendicular}{perpendiculars} from $P$ to the three sides of $\triangle ABC$ and let them meet $\overline{AB}, \overline{BC},$ and $\overline{CA}$ at $D, E,$ and $F$ respectively.  

 
                    \begin{figure}[H]
                        \centering
                        \includegraphics[width=0.5\textwidth]{Solution/AIME_1999_I_14_Solution_Figure_1.png}
                    \end{figure}
                    Let $BE = x, CF = y,$ and $AD = z$.  We have that
 \begin{align*}DP&=z\tan\theta\\ EP&=x\tan\theta\\ FP&=y\tan\theta\end{align*}
We can then use the tool of calculating area in two ways
 \begin{align*}[ABC]&=[PAB]+[PBC]+[PCA]\\ &=\frac{1}{2}(13)(z\tan\theta)+\frac{1}{2}(14)(x\tan\theta)+\frac{1}{2}(15)(y\tan\theta)\\ &=\frac{1}{2}\tan\theta(13z+14x+15y)\end{align*}
On the other hand,
 \begin{align*}[ABC]&=\sqrt{s(s-a)(s-b)(s-c)}\\ &=\sqrt{21\cdot6\cdot7\cdot8}\\ &=84\end{align*}
We still need $13z+14x+15y$ though. We have all these \href{https://artofproblemsolving.com/wiki/index.php?title=Right\_triangle}{right triangles} and we haven't even touched \href{https://artofproblemsolving.com/wiki/index.php?title=Pythagorean\_theorem}{Pythagoras}. So we give it a shot:
 \begin{align}x^2+x^2\tan^2\theta&=z^2\tan^2\theta+(13-z)^2\\ z^2+z^2\tan^2\theta&=y^2\tan^2\theta+(15-y)^2\\ y^2+y^2\tan^2\theta&=x^2\tan^2\theta+(14-x)^2\end{align}Adding $(1) + (2) + (3)$ gives
 \begin{align*}x^2+y^2+z^2&=(14-x)^2+(15-y)^2+(13-z)^2\\ \Rightarrow13z+14x+15y&=295\end{align*}
Recall that we found that $[ABC]=\frac{1}{2}\tan\theta(13z+14x+15y)=84$. Plugging in $13z+14x+15y=295$, we get $\tan\theta=\frac{168}{295}$, giving us $\boxed{463}$ for an answer.

 





\textbf{Solution 2}

Let $AB=c$, $BC=a$, $AC=b$, $PA=x$, $PB=y$, and $PC=z$.

 So by the \href{https://artofproblemsolving.com/wiki/index.php?title=Law\_of\_Cosines}{Law of Cosines}, we have:
 \begin{align*}x^2 &= z^2 + b^2 - 2bz\cos{\theta}\\ y^2 &= x^2 + c^2 - 2cx\cos{\theta}\\ z^2 &= y^2 + a^2 - 2ay\cos{\theta}\end{align*}
Adding these equations and rearranging, we have:
 \[a^2 + b^2 + c^2 = (2bz + 2cx + 2ay)\cos{\theta}\qquad(1)\]Now $[CAP] + [ABP] + [BCP] = [ABC] = \sqrt {(21)(8)(7)(6)} = 84$, by \href{https://artofproblemsolving.com/wiki/index.php?title=Heron\%27s\_formula}{Heron's formula}.

 Now the area of a triangle, $[A] = \frac {mn\sin{\beta}}{2}$, where $m$ and $n$ are sides on either side of an angle, $\beta$. So, 
 \begin{align*}[CAP] &= \frac {bz\sin{\theta}}{2}\\ [ABP] &= \frac {cx\sin{\theta}}{2}\\ [BCP] &= \frac {ay\sin{\theta}}{2}\end{align*}
Adding these equations yields:
 \begin{align*}[ABC]= 84 &= \frac {(bz + cx + ay)\sin{\theta}}{2}\\ \Rightarrow 168&= (bz + cx + ay)\sin{\theta}\qquad (2)\end{align*}
Dividing $(2)$ by $(1)$, we have:
 \begin{align*}\frac {168}{a^2 + b^2 + c^2} &= \frac {(bz + cx + ay)\sin{\theta}}{(2bz + 2cx + 2ay)\cos{\theta}}\\ \Rightarrow \tan{\theta} = \frac {336}{a^2 + b^2 + c^2} &= \frac {336}{14^2 + 15^2 + 13^2} = \frac {336}{590} = \frac {168}{295}\end{align*}
Thus, $m + n = 168 + 295 = \boxed{463}$.

 Note: In fact, this problem is unfairly easy to those who happen to have learned about Brocard point. The Brocard Angle is given by
 \[cot(\theta)=\frac{a^2+b^2+c^2}{4\Delta}\]

 





\textbf{Solution 3}

Let $\angle{PAB} = \angle{PBC} = \angle{PCA} = x.$ Then, using Law of Cosines on the three triangles containing vertex $P,$ we have
 \begin{align*} b^2 &= a^2 + 169 - 26a \cos x \\ c^2 &= b^2 + 196 - 28b \cos x \\ a^2 &= c^2 + 225 - 30c \cos x. \end{align*}
Add the three equations up and rearrange to obtain 
 \[(13a + 14b + 15c) \cos x = 295.\] Also, using $[ABC] = \frac{1}{2}ab \sin \angle C$ we have 
 \[[ABC] = [APB] + [BPC] + [CPA] = \dfrac{\sin x}{2}(13a + 14b + 15c) = 84 \iff (13a + 14b + 15c) \sin x = 168.\] Divide the two equations to obtain $\tan x = \frac{168}{295} \iff \boxed{463}.~\square$

 


 





\textbf{Solution 4 (Law of sines)}

Firstly, denote angles $ABC$, $BCA$, and $CAB$ as $B$, $A$, and $C$ respectively. Let $\angle{PAB}=x$.Notice that by angle chasing that $\angle{BPC}=180-C$ and $\angle{BPA}=180-B$.Using the nice properties of the 13-14-15 triangle, we have $\sin B = \frac{12}{13}$ and $\sin C = \frac{4}{5}$. $\cos C$ is easily computed, so we have $\cos C=\frac{3}{5}$.

 Using Law of Sines,
 \begin{align*} \frac{BP}{\sin x} &= \frac{13}{\sin (180 - B)} \\ \frac{BP}{\sin (C - x)} &= \frac{14}{\sin (180 - C)} \end{align*}
hence, 
 \begin{align*} \frac{BP}{\sin x} &= \frac{13}{\sin B} \\ \frac{BP}{\sin (C - x)} &= \frac{14}{\sin C} \end{align*}
Now, computation carries the rest.
 \begin{align*} \frac{13 \sin x}{\sin B} &= \frac{14 \sin (C-x)}{\sin C}  \\ \frac{169 \sin x}{12} &= \frac{210 \sin (C-x)}{12} \\ 169 \sin x &= 210 (\sin C \cos x - \cos C \sin x) \\ 169 \sin x &= 210 (\frac{4}{5} \cos x - \frac{3}{5} \sin x) \\ 169 \sin x &= 168 \cos x - 126 \sin x \\ 295 \sin x &= 168 \cos x \\ \tan x &= \frac{168}{295} \end{align*}
Extracting yields $168 + 295 = \boxed{463}$.

 





\textbf{Solution 5 (Coordinate Bash)}


                    \begin{figure}[H]
                        \centering
                        \includegraphics[width=0.5\textwidth]{Solution/AIME_1999_I_14_Solution_Figure_1.png}
                    \end{figure}
                    Plot $\triangle ABC$ on the coordinate plane with \(B=(0,0), C=(14,0), A=(5,12).\)Let $\angle PAB=\angle PBC=\angle PCA=a$, and let $P=(x,y)$. For convenience, define $\tan a = k$.

 Since $\angle PBC=a$, the slope of $PB$ is $k$, so \(PB:\ y=kx.\)

 Let $\angle ADB$ be the angle that line $AB$ makes with the $x$-axis. Since $\tan(\angle ADB)=\frac{12}{5}$, the slope of $PA$ is
 \[\tan(\angle ADB+a)=\frac{\frac{12}{5}+k}{1-\frac{12}{5}k}=\frac{12+5k}{5-12k}.\]Thus the equation of $PA$ is
 \[y-12=\frac{12+5k}{5-12k}(x-5),\]which simplifies to
 \[PA:\ y=\frac{12x+5kx-169k}{5-12k}.\]

 Similarly, since $\tan(\angle ECA)=-\frac{4}{3}$, the slope of $PC$ is
 \[\tan(\angle ECA+a)=\frac{-\frac{4}{3}+k}{1+\frac{4}{3}k}=\frac{3k-4}{4k+3}.\]Thus the equation of $PC$ is
 \[PC:\ y=\frac{3k-4}{4k+3}(x-14),\]or equivalently,
 \[PC:\ y=\frac{3kx-4x-42k+56}{4k+3}.\]

 Setting the equations of $PB$ and $PA$ equal gives
 \[\frac{12x+5kx-169k}{5-12k}=kx.\]Clearing denominators yields
 \[12x+5kx-169k=5kx-12k^2x,\]which simplifies to
 \[12x-169k=-12k^2x.\]

 Next, setting the equations of $PB$ and $PC$ equal gives
 \[\frac{3kx-4x-42k+56}{4k+3}=kx.\]Clearing denominators yields
 \[3kx-4x-42k+56=4k^2x+3kx,\]or
 \[-4x-42k+56=4k^2x.\]Multiplying by $-3$ gives
 \[12x+126k-168=-12k^2x.\]

 Equating the right-hand sides of the two derived equations gives
 \[12x+126k-168=12x-169k,\]so \(295k=168,\) and hence \(k=\frac{168}{295}.\)

 Because the problem asks for $m+n$, where $k=\frac{m}{n}$, the final answer is \(168+295=\boxed{463}.\)

 




	\item (\textit{AIME I 1999}) Consider the paper triangle whose vertices are $(0,0), (34,0),$ and $(16,24).$  The vertices of its midpoint triangle are the midpoints of its sides.  A triangular pyramid is formed by folding the triangle along the sides of its midpoint triangle.  What is the volume of this pyramid?


 
 \textbf{Answer: }408



\textbf{Solution 1}


                    \begin{figure}[H]
                        \centering
                        \includegraphics[width=0.5\textwidth]{Solution/AIME_1999_I_15_Solution_Figure_1.png}
                    \end{figure}
                    
                    \begin{figure}[H]
                        \centering
                        \includegraphics[width=0.5\textwidth]{Solution/AIME_1999_I_15_Solution_Figure_1.png}
                    \end{figure}
                    As shown in the image above, let $D$, $E$, and $F$ be the midpoints of $\overline{BC}$, $\overline{CA}$, and $\overline{AB}$, respectively.  Suppose $P$ is the apex of the tetrahedron, and let $O$ be the foot of the altitude from $P$ to $\triangle ABC$.  The crux of this problem is the following lemma.

 \textbf{Lemma:} The point $O$ is the orthocenter of $\triangle ABC$.

 \textit{Proof.} Observe that 
 \[OF^2 - OE^2 = PF^2 - PE^2 = AF^2 - AE^2;\] the first equality follows by the Pythagorean Theorem, while the second follows from $AF = FP$ and $AE = EP$.  Thus, by the Perpendicularity Lemma, $AO$ is perpendicular to $FE$ and hence $BC$.  Analogously, $O$ lies on the $B$-altitude and $C$-altitude of $\triangle ABC$, and so $O$ is, indeed, the orthocenter of $\triangle ABC$.

 To find the coordinates of $O$, we need to find the intersection point of altitudes $BE$ and $AD$. The equation of $BE$ is simply $x=16$. $AD$ is \href{https://artofproblemsolving.com/wiki/index.php?title=Perpendicular}{perpendicular} to line $BC$, so the slope of $AD$ is equal to the negative reciprocal of the slope of $BC$. $BC$ has slope $\frac{24-0}{16-34}=-\frac{4}{3}$, therefore $y=\frac{3}{4} x$. These two lines intersect at $(16,12)$, so that's the base of the height of the tetrahedron. 

 Let $S$ be the foot of altitude $BS$ in $\triangle BPQ$. From the \href{https://artofproblemsolving.com/wiki/index.php?title=Pythagorean\_Theorem}{Pythagorean Theorem}, $h=\sqrt{BS^2-SO^2}$. However, since $S$ and $O$ are, by coincidence, the same point, $SO=0$ and $h=12$.

 The area of the base is $102$, so the volume is $\frac{102*12}{3}=\boxed{408}$.~Shen Kislay Kai

 





\textbf{Solution 2}

Consider the diagram provided in the previous solution. We first note that the medial triangle has coordinates $(17, 0, 0)$, $(8, 12, 0)$, and $(25, 12, 0)$. We can compute the area of this triangle as $102$. Suppose $(x, y, z)$ are the coordinates of the vertex of the resulting pyramid. Call this point $V$. Clearly, the height of the pyramid is $z$. The desired volume is thus $\frac{102z}{3} = 34z$. 

 We note that when folding the triangle to form the pyramid, some side lengths must stay the same. In particular, $VR = RA$, $VP = PB$, and $VQ = QC$. We then use distance formula to find the distances from $V$ to each of the vertices of the medial triangle. We thus arrive at a fairly simple system of equations, yielding $z = 12$. The desired volume is thus $34 \times 12 = \boxed{408}$.~Shen Kislay Kai

 
\textbf{Remark.} It may be easier to instead consider a transformation of the medial triangle to the coordinates $(0,0,0)$, $(17,0,0)$, $(8,0,12)$ (these are simply half of the coordinates of the original triangle, moved to a 3D plane). ~eevee9406

 





\textbf{Solution 3}

The formed tetrahedron has pairwise parallel planar and oppositely equal length ($4\sqrt{13},15,17$) edges and can be inscribed in a parallelepiped (rectangular box) with the six tetrahedral edges as non-intersecting diagonals of the box faces.  Let the edge lengths of the parallelepiped be $p,q,r$ and solve (by Pythagoras)

 $p^2+q^2=4^2\cdot{13}$

 $q^2+r^2=15^2$

 $r^2+p^2=17^2$

 to find that $(p^2,q^2,r^2)=(153,136,72)=(3^2\cdot{17},2^3\cdot{17},2^3\cdot{3^2}).$

 Use the fact that the ratio of volumes between an inscribed tetrahedron and its circumscribing parallelepiped is $\frac{1}{3}$ and then the volume is

 $\frac{1}{3}pqr=\frac{1}{3}\sqrt{2^6\cdot{3^4}\cdot{17^2}}=\boxed{408}$

 
Solution by D. Adrian Tanner

 





\textbf{Solution 4}

Let $A = (0,0), B = (16, 24), C = (34,0).$ Then define $D,E,F$ as the midpoints of $BC, AC, AB$. By Pythagorean theorem, $EF = \frac{1}{2} BC = 15, DE = \frac{1}{2}AB = 4 \sqrt{13}, DF = \frac{1}{2} AC = 17.$ Then let $P$ be the point in space which is the vertex of the tetrahedron with base $DEF$. 

 
Note that $\triangle DEP \cong \triangle EDF$. Create point $F'$ on the plane of $DEF$ such that $\triangle DEP \cong \triangle DEF'$ (i.e by reflecting $F$ over the perpendicular bisector of $DE$). Project $F, P$ onto $DE$ as $X, Y$. Note by the definition of $F'$ then $\angle PYF'$ is the dihedral angle between planes $DEP, DEF$.

 Now see that by Heron's, 
 \[[DEP] = [DEF] = \sqrt{(16 + 2 \sqrt{13})(16 - 2 \sqrt{13})(1 + 2 \sqrt{13})(-1 + 2 \sqrt{13})} = 102.\]So $PY$, the hypotenuse $DEP$ has length $\frac{102 \cdot 2}{4 \sqrt{13}} = \frac{51}{\sqrt{13}}$. Similarly $F'Y = \frac{51}{\sqrt{13}}.$ Further from Pythagoras $DY = \sqrt{DP^2 - PY^2} = \frac{18}{\sqrt{13}}.$ Symmetrically $EX = \frac{18}{\sqrt{13}}.$ Therefore $XY = DE - DY - EX = \frac{16}{\sqrt{13}}.$

 By Law of Cosines on $\triangle PYF'$,
 \begin{align*} PF'^2 &= PY^2 + F'Y^2 - 2 \cdot PY \cdot F'Y \cos{\angle PYF'} \\ PF^2 - XY^2 &= 2 (\frac{51}{\sqrt{13}})^2 \cos{\angle PYF'} \\ (4\sqrt{13})^2 - (\frac{16}{\sqrt{13}})^2 &= 2 (\frac{51}{\sqrt{13}})^2 \cos{\angle PYF'}  \\ \cos{\angle PYF'} &= \frac{9}{17} \\ \sin{\angle PYF'} &= \frac{4 \sqrt{13}}{17}. \end{align*}
.

 Therefore the altitude of the tetrahedron from vertex $P$ to base $DEF$ is $PY \sin{\angle PYF'} = \frac{51}{\sqrt{13}} \frac{4 \sqrt{13}}{17} = 12.$ So the area is $\frac{1}{3}bh = \frac{1}{3} 12 \cdot 102 = \boxed{408}.$ 

 





\textbf{Solution 5 (Formula Abuse)}

The Pyramid is a disphenoid, because opposite sides have the same length. The volume of a disphenoid is given by 

 
 \[V = \sqrt{\frac{(l^{2}+m^{2}-n^{2})(l^{2}-m^{2}+n^{2})(-l^{2}+m^{2}+n^{2})}{72}}.\]

 Using the Pythagorean theorem, the side lengths of the smaller triangle are \(15\), \(4\sqrt{13}\), and \(17\). Plugging in, we get 

 
 \[V = \sqrt{166464} = \boxed{408}.\]

 ~~Disphenoid\_lover

 




\end{enumerate}